
\newpage
\section{Tipps für die Klausur}

Hier sind ein paar Erfahrungen, die ich beim Korrigieren von Klausuren gemacht habe. Aber keine Garantie, solche Dingen ändern sich gerne mal!

\begin{itemize}
	\item \fat{Aufgaben beliebig lösen:} \\		
	Wenn \fat{nicht klar gesagt} wird, wie man eine Aufgabe zu lösen hat, löse sie so, wie du es willst (und vorallem am schnellsten bist). Pass aber auf! Es könnte sein, dass indirekt ein Verfahren deutlich schneller ist und das nicht in der Aufgabe steht!
	
	\item \fat{Logiklücken und Fehler in Aufgaben} \\		
	Wenn du einen Fehler in einer Aufgabe findest oder denkst, dass du einen Weg genommen hast, der nicht bedacht war und der die Aufgabe z.B. deutlich vereinfacht hat, schreibe einen Satz, wie du darauf kommst. Das wird mit einbezogen. Logischerweise gibt es aber keine Punkte, wenn dieser Weg dann falsch ist.
	
	\item \fat{Überprüft eure Lösungen mit Bauchgefühl:} \\
	Es gibt Aufgaben, da kann man aus einer Grafik ablesen, was die Lösung einer Rechenaufgabe (vielleicht auch nur grob) sein müsste. In einer Klausur gab es z.B. eine einfach Quadratur Aufgabe und man konnte aus der Grafik sehen, dass die Fläche unter der Funktion negativ und betragsmäßig recht klein sein musste. Wenn man dann halt auf Zahlen wie $12$ kommt, sollte man sich seine Rechnung vielleicht noch einmal durchgehen. Das gleiche gilt auch für Verständnisaufgaben! Wenn man nämlich (als Beispiel) erkennt, dass seine Lösung falsch ist, aber nicht die Zeit hat, diese zu korrigieren, schreibt als Text hin, dass ihr erkannt habt, dass die Lösung falsch sein muss. Dafür gibt es Punkte, weil wir sehen, dass ihr es verstanden habt.
	
	%In der Klausur wird man bis zu 42 Punkte erhalten können, um die 17 benötigt man zum Bestehen (normale 40 \% Grenze)
	
	\item \fat{Aufbau}:
	\begin{itemize}
		\item In der Klausur gibt es normal 5-6 Abschnitte mit jeweils mehreren Teilaufgaben. Jeder Abschnitt repräsentiert meist ein Oberkapitel. Dementsprechend gibt es mehrere Themen, die nicht abgefragt werden.
		
		\item Es werden \fat{mindestens} 1-2 Aufgaben drankommen, die in der Kategorie der grafischen Aufgaben eingeordnet werden kann. Also Informationen aus einer Grafik auslesen, selber etwas in eine Grafik einzeichnen, anhand einer argumentieren oder oder oder \dots
		
		\item Die letzte Aufgabe der Klausur ist (immer) eine \fat{Programmieraufgabe}. Hierbei gilt, dass man in Java oder Java-ähnlichem Pseudocode geschrieben werden muss. Solange ihr nicht von irgendwelchen Packages ausgeht, die euch alle Funktionalitäten abnehmen und euer Code verständlich ist, alles super!% \\
		%Noch etwas dazu: Die Programmieraufgabe sieht meist doof aus, ist aber billig und gibt für den Zeitaufwand überdurschnittlich viele Punkte. Es lohnt sich also, diese zu bearbeiten.
		
		\item Es gibt immer 2-3 Teilaufgaben, die sehr schwierig sind bzw. \fat{nicht in den Tutorien} behandelt worden sind (und die ich vermutlich ohne Üben auch nicht schaffen würde). Macht euch darüber keine Sorge, ohne diese kann man immer noch eine 1,3 schreiben. \Cooley
	\end{itemize}
	
	\item \fat{Verständnis wichtiger als Rechenkünste:}
	\begin{itemize}
		\item Es gibt mehr Punktabzug, wenn ihr was falsch macht und wir sehen, dass es nicht verstanden wurde. Außerdem gibt es häufiger Verständnisfragen (meist nach einer kleinen Rechenaufgabe).
		\item Wir korrigieren vollständig mit \fat{Folgefehler}. Aus diesem Grund ist Verständnis noch einmal wichtiger als Rechenkünste. Wobei ihr euch natürlich nicht nur auf dem Grundschulniveau der Rechenkünste befinden solltet. Im Allgemeinen gibt jeder mathematische Fehler, der nicht durch Folgefehler entstanden ist, \fat{einen halben Punkt Abzug}.
	\end{itemize}
	
	\item \fat{Textantworten sind ok!}\\
	Wenn ihr Schwierigkeiten habt, probiert eure Antwort wenigstens in Worten zu beschreiben, anstatt es leer zu lassen. Wenn wir anhand von den wenigen Worten erkennen können, dass ihr es eigentlich verstanden habt, \fat{gibt es noch Punkte!}
	
	\item \fat{Die Punktevergabe ist anders je Aufgabenteil.}\\
	Es lohnt sich also, am Anfang der Klausur über alle Aufgaben drüberzuschauen und sich die Aufgaben auszusuchen, die man selber nicht nur schnell lösen kann, sondern die auch mehr Punkte geben.
	
	\item \fat{Zu billige Aufgaben?} \\
	Es kann durchaus seeehr billige Aufgaben in der Klausur geben. Also wenn ihr denkt, dass eine Aufgabe zu einfach um wahr zu sein ist, überprüft nochmal alles. Aber wenn ihr dann weiterhin sie für einfach haltet, ist sie das wahrscheinlich auch.
	
	\item \fat{Hinweise auf dem Klausurblatt} \\
	Bei einigen Aufgaben stehen Hinweise. Die stehen da meist nicht ohne Grund also nutzt sie! Falls ihr aber eine Aufgabe ohne den Hinweis lösen könnt, dann ist das auch super, nur eher unwahrscheinlich. \Smiley
	
	\item \fat{Ruhe} \\
	Das Wichtigste: immer Ruhe bewahren. Numprog ist cool, und so ist auch die Klausur. Macht euch keine Panik und liest euch alles gemütlich durch. Und hey, wenn ihr es nicht besteht, findet Numprog ja immerhin jedes Semester statt. \Winkey
	
	%				Fast garantiert & Hoch & Mittel & Niedrig \\ \midrule
	%				% Inhalt
	%				Gleitkommazahlen 	& LR/QR Zerlegung 				& Power-Iteration 	& Fourier-Transformation \\
	%				Interpolation 		& Fixpunkte/Banach'scher Fixpunktsatz & Jacobi/Gauss-Seidel	& \\
	%				Quadratur 			& Differentialgleichungen 		& Spline Interpolation & \\
	%									& Kondition/Stabilität 			& 					 & \\
	%Und da das auch manche immer (meiner Meinung nach unverständlich) wissen wollen, im Durschnitt beträgt die \fat{Durchfallquote ca. 15-25 \%}, ein Semester war sie mit \fat{41 \%} erstaunlich hoch. 
\end{itemize}

%----------------------------
\newpage
\section{Typische Fehler}

Hier liste ich mal ein paar Fehler auf, die mir persönlich beim Korrigieren von Numprog Klausuren aufgefallen sind. Alles etwas querbeet und chaotisch, aber vielleicht trotzdem hilfreich. \Smiley

\begin{itemize}
	\item Aufgabenstellung nicht sorgfältig genug durchgelesen und Teilaufgaben deswegen nicht beantwortet. 
	\begin{itemize}
		\item Gerade so Teile wie "allgemeine Formeln" angeben werden häufiger vernachlässigt. Darauf gibt es Punkte, sonst würde es nicht in der Aufgabe stehen!
		\item Wenn nach einer "kleinsten Zahl größer Null" gefragt ist, darf die Antwort keine negative Zahl sein.
	\end{itemize}
	\item Hinweise einer Aufgabe nicht durchgelesen. Zum Teil stehen da alle Formeln, die man für die Aufgabe benötigt.
	\item Wenn eine Zahlendarstellung $x$ Bytes zur Verfügung hat, ist das nicht äquivalent zu $x$ Bits, sondern $8 \cdot x$ Bits.
	\item Umwandlung zu IEEE: 
	\begin{itemize}
		\item Null steht beim Vorzeichenbit für eine positive Zahl!
		\item Bei der Bestimmung des Exponenten vergessen, den konstanten Offset (die $-127$) einzuberechnen.
		\item Die Anzahl an Mantissenbits vertan (es sind 23 Bits, ohne die erste Eins, da die nicht abgespeichert werden muss!) \& vergessen, dass man runden muss (falls die Mantisse eigentlich länger als 23 Bits ist)
		\item Klar spezifizieren, welche Teile Vorzeichen, Exponent und Mantisse ist.
	\end{itemize}
	\item Die Ableitung von $x^2$ ist $2x$ \dots
	\item Stabilität \& Kondition verwechselt. Das ist nicht dasselbe!
	\item Für eine Limes Untersuchung kann häufig der "L'Hopital" weiterhelfen.
	\item Bei der LR-Zerlegung nach der eigentlichen Zerlegung und Vorwärtssubstitution aufgehört, obwohl noch die Rückwärtssubstitution durchgeführt werden muss.
	\item Beim Lösen eines überbestimmten Gleichungssystem einfach am Anfang die Matrix so ausschneiden, dass sie quadratisch wird, ist nicht erlaubt!
	\item Bei einem überbestimmten Gleichungsystem zum Lösen Gauß-Eliminationsverfahren anwenden ist Quatsch.
	\item Polynome sind stets beliebig häufig stetig differenzierbar.
	\item Wenn bei einer Quadratur Rechenaufgabe etwas negatives herauskommt und man bei der Zeichnung eine Fläche einzeichnet, die insgesamt positiv ist sollte man vielleicht noch einmal über seine Rechnung drüberschauen.
	\item Bei Lagrange-Basen den Aufwand machen, das LGS aufzustellen, obwohl die Matrix mit Lagrange eh immer einer Einheitsmatrix entspricht.
	\item Implizites Euler-Verfahren \fat{ist NICHT} genauer als Explizites Euler-Verfahren. Beide haben Ordnung 1. Das implizite ist stabiler, aber nicht (allgemein) genauer!
	\item Runge Kutta ist mit 4. Ordnung genauer als Heun mit 2. Ordnung, beide sind genauer als Explizites- sowie Implizites Euler Verfahren, da diese nur Ordnung 1 haben.
	\item Bei einer Gauß-Quadratur Aufgabe die Gewichtungen ausrechnen, obwohl diese in der Aufgabe gegeben sind.
	\item Zahlen aus der Aufgabe falsch abgeschrieben\dots
	\item Wenn eine Iterationsformel ala $x_{k+1} = x_k + a(x_k)$ gegeben ist, ist die Iterationsvorschrift einfach nur $\Phi(x) = x + a(x)$.
	\item Einzeichnen einer Trapezregel: es heißt nicht umsonst Trapezregel, also eine Gerade oder Fünfeck einzeichen ist irgendwo falsch. Es kann aber durchaus sein, dass ein Dreieck oder Rechteck entsteht.
	\item Bei grafischen Aufgaben sich die Koordinatenachsen nicht richtig angeschaut. Manchmal ist links unten nicht der Ursprung. Also schaut erst einmal an, wo $(0,0)$ überhaupt ist.
	\item Wenn eine Frage etwas ist wie "Was beobachten Sie? Warum?", dann ist das eigentlich nie zum Spaß da. Falls es also nichts zu beobachten gibt hast du vielleicht etwas falsch ausgerechnet.
	\item Nicht in der Lage sein, die Winkelhalbierende einzuzeichnen (das ist eine Gerade mit Steigung 1, die durch den Ursprung geht)
\end{itemize}

